% !TEX program = xelatex
\documentclass[10pt]{article}
\usepackage{graphicx}
\usepackage{float}
\usepackage{listings}
\usepackage{xcolor}
\usepackage{caption}
\usepackage{algorithm}
\usepackage{amsmath}
\usepackage{amssymb}
\usepackage{algpseudocode}
\usepackage{fontspec} % Crucial for XeLaTeX to use system fonts
\usepackage{xunicode} % For extra character support
\usepackage{xltxtra}  % Extra fixes for XeLaTeX
\usepackage[top=2.5cm, bottom=2.5cm, left=3cm, right=2.5cm]{geometry}
\usepackage[hidelinks]{hyperref}
\usepackage{booktabs}

% Set the main font (You can change these to any font on your system)
\setmainfont{QTBookmann} % A good, free Times New Roman alternative
\setmonofont{VictorMono Nerd Font}

% Title Page Information
\title{A Simple Report on Various Topics}
\author{Your Name}
\date{\today} % Automatically inserts the current date

\definecolor{codegreen}{rgb}{0,0.6,0}
\definecolor{codegray}{rgb}{0.5,0.5,0.5}
\definecolor{codepurple}{rgb}{0.58,0,0.82}
\definecolor{backcolour}{rgb}{0.95,0.95,0.92}

\lstdefinestyle{mystyle}{
    backgroundcolor=\color{backcolour},   
    commentstyle=\color{codegreen},
    keywordstyle=\color{magenta},
    numberstyle=\tiny\color{codegray},
    stringstyle=\color{codepurple},
    basicstyle=\ttfamily\footnotesize,
    breakatwhitespace=false,         
    breaklines=true,                 
    captionpos=b,                    
    keepspaces=true,                 
    numbers=left,                    
    numbersep=5pt,                  
    showspaces=false,                
    showstringspaces=false,
    showtabs=false,                  
    tabsize=2,
    frame=single
}

\lstset{style=mystyle}

\title{Shrinking and Stretching of Images}
\author{Brandon Marquez Salazar}
\date{}

\begin{document}

\maketitle

\section{Objective and Approach}

In the current work, the objective was to modify an existing algorithm to achieve independent shrinking and stretching of a given image along its width and height. The starting point was the in-class program 07, which uses a single scaling factor $f$. This original algorithm handles two cases: $f > 1$ for enlargement and $f < 1$ for reduction.

The original algorithm is as follows:

\begin{algorithm}
\caption{Image Scaling (Escalar)}
\begin{algorithmic}[1]
\Procedure{Escalar}{$I_{in}$, $I_{out}$, $f$}
    \If{$f > 1.0$} \Comment{Upscaling: nearest neighbor}
        \State $s \gets f$, $I_{out}.\text{create}(h \cdot s, w \cdot s)$
        \For{$u \in [0,h)$}
            \For{$v \in [0,w)$}
                \State $p \gets I_{in}[u,v]$
                \For{$i \in [u \cdot s, (u+1) \cdot s)$}
                    \For{$j \in [v \cdot s, (v+1) \cdot s)$}
                        \State $I_{out}[i,j] \gets p$
                    \EndFor
                \EndFor
            \EndFor
        \EndFor
        
    \Else \Comment{Downscaling: average pooling}
        \State $s \gets \lfloor 1/f \rfloor$
        \If{$h \not\equiv 0 \mod s$ OR $w \not\equiv 0 \mod s$}
            \State $I_{out} \gets I_{in}$ \Comment{Dimension error}
        \Else
            \State $I_{out}.\text{create}(h/s, w/s)$
            \For{$u \in [0,h)$ step $s$}
                \For{$v \in [0,w)$ step $s$}
                    \State $I_{out}[u/s,v/s] \gets \frac{1}{s^2} \sum_{i=u}^{u+s-1} \sum_{j=v}^{v+s-1} I_{in}[i,j]$
                \EndFor
            \EndFor
            \State \Call{AjusteDeRango}{$I_{out}$}
        \EndIf
    \EndIf
    \State $I_{out}.\text{convertTo}(CV\_8UC1)$
\EndProcedure
\end{algorithmic}
\end{algorithm}

\section{Proposed Algorithm}

The proposed new algorithm expands the original logic to accommodate two independent scaling factors, $f_x$ and $f_y$, for the x-axis (width) and y-axis (height), respectively.

The proposed algorithm is:

\begin{algorithm}
\caption{Asymmetric Image Resize}
\begin{algorithmic}[1]
\Procedure{Resize}{$f_x$, $f_y$}
    \State $I_{in} \gets \text{this} \rightarrow \text{image}$ \Comment{Input image $I_{in}[u,v]$}
    
    \State $(s_x, s_y) \gets \begin{cases}
        (f_x, f_y) & \text{if } f_x > 1, f_y > 1 \\
        (\lfloor 1/f_x \rfloor, f_y) & \text{if } f_x < 1, f_y > 1 \\
        (f_x, \lfloor 1/f_y \rfloor) & \text{if } f_x > 1, f_y < 1 \\
        (\lfloor 1/f_x \rfloor, \lfloor 1/f_y \rfloor) & \text{otherwise}
    \end{cases}$
    
    \State \textbf{Output dimensions:} $(h_{out}, w_{out}) \gets \begin{cases}
        (h \cdot s_y, w \cdot s_x) & \text{if } f_x > 1, f_y > 1 \\
        (h \cdot s_y, w / s_x) & \text{if } f_x < 1, f_y > 1 \\
        (h / s_y, w \cdot s_x) & \text{if } f_x > 1, f_y < 1 \\
        (h / s_y, w / s_x) & \text{otherwise}
    \end{cases}$
    
    \If{$\exists$ downscaling AND $(w \not\equiv 0 \mod s_x$ OR $h \not\equiv 0 \mod s_y)$}
        \State \Return error
    \EndIf
    
    \For{$(i,j) \in [0,h_{out}) \times [0,w_{out})$}
        \State $(u,v) \gets \Call{MapCoordinates}{i,j,s_x,s_y,f_x,f_y}$
        \State $I_{out}[i,j] \gets \Call{ComputePixel}{I_{in},u,v,s_x,s_y,f_x,f_y}$
    \EndFor
    
    \If{$f_x < 1$ AND $f_y < 1$} \Call{AjusteDeRango}{$I_{out}$} \EndIf
    \State $\text{this} \rightarrow \text{image} \gets I_{out}$
\EndProcedure
\end{algorithmic}
\end{algorithm}

\begin{algorithm}
\caption{Pixel Coordinate Mapping}
\begin{algorithmic}[1]
\Function{MapCoordinates}{$i$, $j$, $s_x$, $s_y$, $f_x$, $f_y$}
    \If{$f_x > 1$ AND $f_y > 1$} \Return $(i/s_y, j/s_x)$ \EndIf
    \If{$f_x < 1$ AND $f_y > 1$} \Return $(i/s_y, j \cdot s_x)$ \EndIf  
    \If{$f_x > 1$ AND $f_y < 1$} \Return $(i \cdot s_y, j/s_x)$ \EndIf
    \If{$f_x < 1$ AND $f_y < 1$} \Return $(i \cdot s_y, j \cdot s_x)$ \EndIf
\EndFunction
\end{algorithmic}
\end{algorithm}

\begin{algorithm}
\caption{Pixel Computation}
\begin{algorithmic}[1]
\Function{ComputePixel}{$I_{in}$, $u$, $v$, $s_x$, $s_y$, $f_x$, $f_y$}
    \If{$f_x > 1$ AND $f_y > 1$} 
        \State \Return $I_{in}[\lfloor u \rfloor, \lfloor v \rfloor]$ \Comment{Nearest neighbor}
    \EndIf
    
    \If{$f_x < 1$ AND $f_y > 1$}
        \State $sum \gets \sum_{y=u}^{u+s_y-1} \sum_{x=v}^{v+s_x-1} I_{in}[y,x]$
        \State \Return $sum / (s_x \cdot s_y)$ \Comment{Average pooling in x}
    \EndIf
    
    \If{$f_x > 1$ AND $f_y < 1$}
        \State $sum \gets \sum_{y=u}^{u+s_y-1} \sum_{x=v}^{v+s_x-1} I_{in}[y,x]$
        \State \Return $sum / (s_x \cdot s_y)$ \Comment{Average pooling in y}
    \EndIf
    
    \If{$f_x < 1$ AND $f_y < 1$}
        \State $sum \gets \sum_{y=u}^{u+s_y-1} \sum_{x=v}^{v+s_x-1} I_{in}[y,x]$
        \State \Return $sum / (s_x \cdot s_y)$ \Comment{Average pooling both}
    \EndIf
\EndFunction
\end{algorithmic}
\end{algorithm}


This expansion from one to four primary cases (combinations of $f_x \lessgtr 1$ and $f_y \lessgtr 1$) introduced significant complexity to the original program structure. The core challenge was managing the independent scaling along each dimension.

\section{Results and Implementation Issues}

Unfortunately, the implemented code demonstrated several problems in its output. These artifacts are likely attributable to the construction and pixel-access logic of the new image matrix, which proved highly sensitive to the independent handling of the two scaling factors.

The program output is as follows:
\begin{lstlisting}[backgroundcolor=\color{gray!10}, basicstyle=\ttfamily\small, frame=single]
g++ code/tareas/tarea02/main.cxx  code/src/ProcesamientoDigitalImagenes.cxx   `pkg-config --cflags --libs opencv4` -lm -g -o code/tareas/tarea02/exe
cd code/tareas/tarea02 && ./exe
Nombre:             ../../res/testImg/C01.bmp
Número de filas:    480
Número de columnas: 640
Total de pixeles:   307200
Número de canales:  1
Intensidad máxima:  nan
Intensidad mínima:  nan
Nombre:             ../../res/testImg/C02.bmp
Número de filas:    0
Número de columnas: 0
Total de pixeles:   0
Número de canales:  1
Intensidad máxima:  nan
Intensidad mínima:  nan
Nombre:             ../../res/testImg/Family.png
Número de filas:    1200
Número de columnas: 1600
Total de pixeles:   1920000
Número de canales:  4
Intensidad máxima:  nan
Intensidad mínima:  nan
terminate called after throwing an instance of 'cv::Exception'
  what():  OpenCV(4.12.0) ../modules/highgui/src/window.cpp:973: error: (-215:Assertion failed) size.width>0 && size.height>0 in function 'imshow'

Aborted (core dumped)
make: *** [Makefile:52: tarea02] Error 134
\end{lstlisting}

\section{Conclusion and Future Work}

Due to a lack of time, the resolution of these errors is deferred to a future effort. The primary conclusion of this work is that the complexity of this two-factor scaling approach requires meticulous attention to the image matrix creation and pixel access patterns. Future work will focus on correcting the identified implementation flaws.

\end{document}
